\chapter{Two or more Inhibiting Systems}

Many examples were found where there were two or possibly more sets of
inhibiting systems, and each set required different conditions for destruction. Eranthis
hyemalis is an excellent example of this and was studied in detail. Eranthis has to be
one of the great horticultural treats. In late March great drifts of the yellow buttercups
carpet the woods and field, and the air is filled with their intense honey fragrance. On
warm days the honey bees and butterflies that have hibernated come for their first
feast after the long winter. As a result there is always a rich crop of seed that ripens in
the first week of May.

Eranthis hyemalis has one inhibiting system that is destroyed at 70 and another
that is destroyed at 40, and they must be destroyed in that order. The seeds are placed
in moist towels at 70 as soon as they are ripe in May. After 3 m at 70, the seeds are
shifted to 40. After an induction period of exactly 55 days (note the length),
germination begins and proceeds at a rate of 250/old (note the high rate) so that
germination is over 90\% complete in tour days. Over the following month the radicle
develops and the cotyledons start to open. The seedlings have to grow under
relatively cool conditions as this species is programmed to go dormant if temperatures
above 70 persist. As a result shifting the seedlings directly to 70 leads to their dieing
off. Of course the simplest procedure for a gardener is to broadcast the seed as soon
as collected and in a few years large colonies will develop.

If the seed is subjected to a three cycle pattern starting at 40, the results are 40-
70-40(90\% in 55-59 d). The first three months at 40 have absolutely no effect on the
percent germination, the induction period, or the rate of germination and is truly a
dormant state with little or no metabolic activity. The-three months at 70 in either the
70-40 or 40-70-40 pattern is again a time of metabolic activity inside the seed despite
the absence of any outward signs, and chemical reactions are occurring that are
absolutely essential for the germination in the following cycle at 40. -
The importance of low temperature metabolism has been a dominant theme in
this book. In Eranthis hyemalis the germination, development of radicle, and
expansion of cotyledons takes place rapidly at 40 and will not occur at 70. It is evident
that in the plant kingdom neither germination inhibitor destruction nor growth of roots
and leaves are restricted to either high (70) or low (40) temperatures.

A number of experiments were conducted on the seeds of Eranthis hyemalis to
better understand the chemical reactions and metabolism in the three months at 70. If
the period at 70 is shortened from three months to two months, germination at 40 is
seriously affected. The induction time lengthens to 60-70 days and the germination
rate drastically slows down. It is evident that two months is not long enough for the
requisite chemical changes to be completed. If the time at 70 is lengthened from three
months to six months, the germinatiion behavior at 40 is identical to seeds given three
months at 70 showing that the additional three months were a period of inactivity.
However as the time at 70 is further increased, the induction period starts to lengthen
and the germination rate declines showing that irreversible injury and death is
beginning to occur. Although it was not examined, it is likely that seedlings from non-
optimum conditions for germination will be weaker and perhaps fail to develop. In the
40 .70-40 treatment the effect of lengthening the first three months period at 40 was not
examined. Presumably it could be lengthened to apoint where irreversible injury
appeared.

Seed of Eranthis hyemalis will not tolerate dry storage for six months at either
70 or 40 and seed so treated is dead. In the future if seeds of Eranthis are to be stored
and distributed, it is recommended that they be stored in moist towels at 70 or 40 and
distributed in this state. As it stands now, all the seed distributed commercially or by
seed exchanges is DOD (dead on delivery).

Eranthis hyemalis produces only the pair of cotyledons in the first year of
growth. Among dicotyledons this behavior is uncommon, but it has been found in
Oxytropis williamsii in Fabaceae; Anemonastrum, Clematis albicoma, and Eranthis in
Ranunculaceae; Dodecatheon amethystinum and Dodecatheon media in
Primulaceae; and Skimmia japonica in Rutaceae.

Suppose that in Eranthis hyemalis the germination had been fast at 70 and zero
at 40. The seed would then have required three months at 70, three months at 40, and
finally germination after a shift to 70. Now Eranthis does not do this, but it illustrates
that the inhibitor destruction pattern for 70-40 germinators can be the same as for 70-
40-70 germinators, and this introduces us to the world of multicycle germinators.
Many species did not germinate until after two or more alternating three month
cycles. These are termed multicycle germinators. While more work is needed to
conclusively demonstrate that two or more sets of inhibitors were present, this is the
most likely interpretation (although it is not readily distinguished from situations where
the rate of germination is very slow at the temperature of the second cycle). These
multicycle germinators were of two types. In the simpler type germination was
reasonably confined to three or four cycles. Actaea rubra and Iris graeberiana were
40-70-40 germinator and Cercis chinensis and Erythronium mesochorum were largely
70-40-70 germinators. Delphinium tricorne was a 70-40-70-40 germinator. In others
the germination dragged on cycle after cycle, usually with a little germination in the first
two cycles with the majority in subsequent cycles.

There is always the question as to whether longer cycle times, different
temperatures, or some other untried factor would have speeded up the process. In fact
some species that were initially observed to have long drawn out germination over
many cycles were later found to germinate readily in outdoor treatment, in various
phototreatments, or when treated with gibberelins. Some Iris have these long drawn
out germinations, particularly the Juno and Oncocyclus. Attempts were made to speed
up germination by subjecting the seeds to month long soakings in water, by removing
the brown seed coats, by light, and by treatment with gibberelin, all without effect.
Combinations of wet and dry cycles and temperatures over 100 were not tried, and
these merit exploration since these species come from regions with hot dessicating
summers.

Table 7-1 summarizes the data on some three cycle germinators. Table 7-2 lists
some species with even longer drawn out germinations. Sometimes these long drawn
out germinations were associatedwith stepwise germination in which the radicle
developed in one cycle and the leaves in a following cycle as in Jeffersonia diphylla
and dubia. One obvious survival advantage of such long drawn out germinations is
that if a particular year had inclement weather, there would be new seedlings for the
following years. However, the way Oncocyclus Iris seed will remain in a firm and
healthy condition for years suggests that some key factor needed for germination is as
yet undiscovered.

\begin{table}[h]
\refstepcounter{table}
\begin{center}
Table \thetable. Species Whose Germination Is Predominantly in the Third or Fourth Cycle
\end{center}

\par
\begin{tabular}{|l|l|}
\hline

\multicolumn{2}{|c|}{70-40-70 Germination } \\
\hline
Acer spicatum & 70-40-70(28\% in 3-8 weeks) \\
 & 40-70-40-70(9\%) \\
Chionanthus virginica & 70-40-70(54\% in 5-9 weeks)-40-70(29\% in 5-7 weeks) \\
 & 40-70(7\%)-40-70(54\% in 4-8 weeks) \\
Corylus cornuta & 70-40-70(86\% in 1-3 weeks) \\
 & 40-70(14\%)-40-70(57\% in 1-6 weeks) \\
Dicentra spectabilis & 70-40(12\%)-70(28\% in 3-9 days) \\
 & 40(21\%)-70(5\%)-40(12\%)-70(1 \%) \\
Dionysia involucrata & 70-40-70(50\%) \\
Iris hookeri & 70-40-70(100\% in 2nd week) \\
Lilium chalcedonicum & 70-40-70(100\% in 5-7 weeks) \\
Paris quadrifolia & 70(6\%)-40-70(69\% in 7-12 weeks) \\
 & 40-70-40-70(66\% in 7th week) \\
Passiflora edulis & 70-40-70(81\% in 4-6 weeks) \\
 & 40-70-40-70 none \\
\hline
\multicolumn{2}{|c|}{ 40-70-40 Germination } \\
\hline
Iris graeberiana      & 40-70-40(90\% in 6-11 weeks) \\
 & 70-40(45\% in 10th week)-70-40(18\% in 9-12 weeks) \\
\hline
\multicolumn{2}{|c|}{ 40-70-40-70 Germination } \\
\hline
Stylophorum diphyllum & 40-70-40(4\%)-70(69\% in 1-15 days) \\
\hline
\multicolumn{2}{|c|}{ 70-40-70-40 Germination } \\
\hline
Viburnum dilitatum    & 70-40-70-40(70\% in 2-5 weeks) \\
Viburnum opulis       & 70-40(3\%)-70-40(62\% in 2-8 weeks) \\
\hline

\end{tabular}
\label{table:s34}
\end{table}

\begin{table}[h]
\refstepcounter{table}
\begin{center}
Table \thetable. Examples of Extended Multicycle Germinations
\end{center}
\par
\begin{tabular}{|l|l|}
\hline
Allium karataviense & 70-40-70(6\%)-40(30\%)-70(10\%)-40(14\%)-70(6\%) \\
                    & 40-70-40(3\%)-70(14\%)-40(2\%)-70(14\%) \\
Eleaganus umbellata & 70-40(24\%)-70(25\%)-40(6\%)-70(18\%) \\
Fritillaria pallida & 40-70-40-70(20\%)-40(60\%) \\
Iris acutiloba & 70-40-70(10\%)40-70-40-70(70\%) \\
Jeffersonia dubia & 70-40(4\%)-70-40(2\%)-70(4\%)-40(14\%)-70(24\%)-40(6\%) \\
Rhodotypos tetrapetala & 70-40(44\%)-70(16\%)-40(12\%) \\
                       & 40(14\%)-70(24\%)-40-70(14\%)-40(11\%)-70(3\%) \\
Taxus baccata       & 70-40-70(1\%)-40-70(28\% in 5-15 days)-40-70 \\
                    & 40-70-40-70(11\% in 1-16 days)-40-70(26\%)-40-70(10\%)-40-70(8\%) \\
\hline
\end{tabular}
\label{table:smc}
\end{table}

